% Project Summary/Abstract (SF 424 Project Abstract; upload for R&R Other Project Information, Field 7)
% ONR FY27 YIP, NOFO N0001426SF004. Compile with pdfLaTeX. The NOFO limit is one page or
% 4,000 characters including spaces, whichever is less, in standard keyboard characters.
% The text is the same as in project-summary.md (the labels are bold here); edit both together.
% Apostrophes are \textquotesingle{} and ligatures are off, so text copied from the PDF stays ASCII.
\documentclass[12pt,letterpaper]{article}
\usepackage[letterpaper,margin=1in]{geometry}
\usepackage[T1]{fontenc}
\usepackage{textcomp}
\usepackage{mathptmx}
\usepackage{microtype}
\DisableLigatures{encoding = *, family = *}
\usepackage[hidelinks]{hyperref}
\hypersetup{pdftitle={Project Summary/Abstract: Trustworthy Collective Reasoning and Evaluation for AI Agents},pdfauthor={Xinya Du}}
\setlength{\parindent}{1.5em}
\setlength{\parskip}{0pt}
\hyphenpenalty=10000 \exhyphenpenalty=10000
\pagestyle{empty}
\begin{document}
{\centering
{\large\bfseries Trustworthy Collective Reasoning and Evaluation for AI Agents\par}
\vspace{3pt}
{\itshape Approved for Public Release}\par}
\vspace{3pt}
\noindent AI agents built on large language models increasingly evaluate one another and decide together: they judge answers, review papers, supply the rewards that train new models, and vote or discuss. Yet a vote counts a mistake that several agents share as independent evidence, dishonest agents can collude, and AI judges rule on text alone, even when one database query would settle a dispute. Agents trained to agree with the majority can learn to agree on mistakes. If five agents tell a Navy watch officer whether a cargo vessel has left port, and three say yes, two because they read the same outdated schedule and one because it is compromised, the vote is wrong: the ship never left port, and an officer who trusts the vote wastes a patrol searching for it at sea.

This project will enable AI agents of unknown and correlated reliability, some possibly dishonest, to reach trustworthy collective decisions and evaluations through debate. In a debate, two agents argue for opposite answers until the dispute narrows to one step an impartial judge can verify.

Research goal 1 (RG1) weighs answers with a statistical reliability model, fitted without an answer key, in which a group of agents sharing a mistake counts at most once; when the model is unsure, even of a unanimous answer, the question goes to debate. RG1 then chooses each side\textquotesingle{}s debater from peer ratings that no agent can use to get itself chosen, and draws the judge by lottery. RG2 builds, to our knowledge, the first agentic harness for peer evaluation by AI: debaters back their steps with tool calls such as database queries, a control loop stops the debate at the steps in dispute, and the judge re-runs the calls before it decides. Planted cases with known answers and random human audits catch colluding debaters and judges. RG3 extends reinforcement learning with verifiable rewards beyond questions with an answer key, such as math and code, to the agents\textquotesingle{} own collective decisions, without human labels. It learns mainly from decisions the judge settled by re-running tool evidence, with rewards that agreeing with one\textquotesingle{}s own group cannot raise. We test the methods on math and code, multi-document questions with maritime scenarios, and scientific peer review, against a red team that trains agents to collude.

If successful, the project will show, in theory and experiments, when agreement among AI agents is evidence, what an impartial judge that re-runs evidence guarantees, and how few audits keep agents honest. It will release an open debate harness, audit protocols, and benchmarks. Because people audit only a small share of decisions and RG3\textquotesingle{}s rewards need no human labels, the methods can scale to AI whose answers people cannot easily evaluate.

The Department of War is developing AI agents for battle management and decision support, and ONR will explore validating AI results by independent AI critique. Agents built on similar models share mistakes, and an adversary need only compromise one; our AI critique withstands both, and each decision states its reasons. Marine Corps guidance says to verify all generative AI outputs before use, but people cannot check them all. The methods serve Division 311\textquotesingle{}s focus on rapid, accurate decision making in networked environments, wherever staff need one decision from several AI tools whose reports may be copied, outdated, forged, or spoofed: maritime domain awareness, intelligence analysis, command and control, and testing AI before fielding. NRL and NIWC Pacific could evaluate them on Navy data, and Division 311\textquotesingle{}s Command and Control program could build them into decision aids.
\end{document}
